\section{Data and Simulation}
\label{sec:data}

\textbf{Real cars: LiRA-CD.}
LiRA-CD~\cite{skar2023lira,lira_data} contains GreenMobility Renault Zoe cars instrumented with an AutoPi telematics unit (three-axis accelerometer, GPS) connected to the CAN bus, and surveys of the same lanes by the Danish Road Directorate's P79 profilometer.
We use the four routes with car data: two urban roads in Copenhagen (CPH1, part of the O2 ring road, and CPH6) and two motorways (M3, M13), each in both directions.
The P79 files provide the wheel-path laser elevations at 0.1\,m with GPS.
In the car files the accelerometer is logged at a nominal \liraRateLow\,Hz (\liraRateLowShare\,\% of passes) or \liraRateHigh\,Hz, quantised to 0.01\,g in most passes, with a median mounting tilt of \liraTiltMedian$^\circ$; speed comes from the CAN bus and position from GPS at \liraGpsHzMin--\liraGpsHzMax\,Hz.
After map matching (Section~\ref{sec:method}) the data contain \liraUnits{} segment-passes over \liraSegments{} distinct 100\,m segments (\liraKm\,km) from \liraPasses{} passes of at least \liraDevices{} cars (telematics units).
Urban segments have mean IRI \liraIriUrban\,m/km and are driven at \liraSpeedUrban\,km/h on average, motorway segments \liraIriMotorway\,m/km at \liraSpeedMotorway\,km/h.
The same cars drive several routes, so leave-one-route-out evaluation is route-disjoint but not vehicle-disjoint for estimators trained on real data (Section~\ref{sec:res_real} tests the consequence).

\textbf{Labels.}
Labels are the IRI per 100\,m, mean of both wheel paths, recomputed from the P79 profiles with an implementation of the quarter-car algorithm~\cite{sayers1995iri,astm_e1926} that reproduces the P79's own IRI ($r=$\,\iriValR{} over \iriValN{} segments; an algorithm check, since both come from the same profiles).
Labels derived from authored game-engine geometry are not equivalent: the curvature-based mesh correction of our earlier pipeline lowers the IRI of these real urban profiles by \legacyCorrUrbanMape\,\% (Supplement~\ref{S-sec:s_earlier}).

\textbf{Parametric vehicle simulator.}
A linear four-degree-of-freedom half-car (bounce, pitch, front and rear unsprung masses) is driven by a road profile under the front axle and the same profile delayed by the wheelbase under the rear axle, after a 250\,mm tyre-envelope moving average; on measured profiles the input is the mean of both wheel paths (no roll), and simulated units are labelled exactly like the real ones.
The solver deviates by at most \simRkDev{} (relative) from a fine-step Runge--Kutta integration of the continuous system.
Vehicle parameters are drawn from ranges for passenger cars~\cite{gillespie1992} (sprung mass 1000--1800\,kg, wheelbase 2.4--3.0\,m, body frequency 1.0--1.6\,Hz, damping ratio 0.20--0.45; Table~\ref{S-tab:simranges}), speed independently of roughness (20--110\,km/h with fluctuations and braking events).
The sensor model applies an optional holder resonance, an anti-aliasing filter, a mounting rotation, irregular sampling, scale and bias errors, white noise, quantisation and sample loss.
Simulated and real signals then go through identical preprocessing code.

\textbf{Simulator variants.}
The reference simulator \textbf{V0} was fixed before any real test result: 30 random vehicles, measured P79 profiles of the training routes, and a 50/50 mixture of a telematics-like sensor (55--62\,Hz, 0.01\,g, rigid mount) and a phone-like sensor (90--210\,Hz, 0.6\,mg) with a holder resonance at 15--40\,Hz.
To find out which modelling choices matter, we compare it with ten variants, listed in a protocol written after an independent review and before the analyses reported here were run (the zero-shot results of V3--V6 were known from our earlier analysis and those of V1 and V2 from the review, so this comparison is diagnostic, not confirmatory): the telematics-like half alone (\textbf{V1}) and the phone-like half alone (\textbf{V2}) of V0; the two missing cells of a 2$\times$2 design, telematics-like with the holder resonance (\textbf{V7}) and phone-like without it (\textbf{V8}); an ideal sensor (\textbf{V3}); a single telematics sensor matched to statistics measured on the unlabelled LiRA-CD streams (\textbf{V5}) and its 50/50 mixture with the phone-like sensor (\textbf{V6}); one mid-range vehicle instead of 30 (\textbf{V4}); and synthetic ISO-8608-type profiles~\cite{iso8608} alone (\textbf{V9}) or added to the measured ones (\textbf{V10}).
The simulated drives never use the profile of the test route.

\textbf{Earlier BeamNG corpus.}
As a negative reference we also train on the corpus of our earlier study: \bngTrips{} manually driven trips of \bngCars{} BeamNG.tech vehicles on \bngMaps{} maps, with labels from ray-cast elevation; its acquisition duplicated \bngImuDupMin--\bngImuDupMax\,\% of IMU samples (Supplement~\ref{S-sec:s_earlier}).
